Esta seção apresenta os principais achados do estudo e sua interpretação à luz
da fundamentação teórica. Os resultados devem ser expostos de forma clara,
utilizando tabelas, gráficos ou descrições conforme necessários. Na discussão,
esses achados são analisados e comparados com estudos anteriores, destacando
suas implicações, limitações e possíveis explicações. Essa seção deve
evidenciar se os objetivos da pesquisa foram alcançados e contribuir para o
entendimento do tema investigado.

\nota{\textbf{Para o Relatório Parcial (resultados parciais):} se a pesquisa
avançou a ponto de ter alguns resultados, apresente-os aqui. Faça uma breve
discussão sobre o que foi alcançado até o momento. Utilize tabelas, gráficos,
fotos, figuras e ilustrações que facilitem a apresentação ordenada dos
resultados.}

\section{Xxxxxx}

Abaixo seguem um exemplo de subseção e figura com suas formatações. Em seguida
apresenta-se um exemplo de tabela (figura~\ref{fig:exemplo} e
tabela~\ref{tab:exemplo}). Troque o conteúdo de exemplo pelos seus resultados.

\begin{figure}[H]
  \caption{Ponto de colheita do gladíolo no estágio R2 da Escala PhenoGlad de desenvolvimento.}
  \label{fig:exemplo}
  \centering
  \begin{minipage}{320pt}
    \centering
    \fbox{\parbox[c][5cm][c]{300pt}{\centering (Substitua esta caixa pela figura)}}\par
    \fonte[4pt]{Adaptada de xxxxx et al. (2019).}
  \end{minipage}
\end{figure}

\begin{table}[H]
  \caption{xxxxxxxxx}
  \label{tab:exemplo}
  \tabelaibge\raggedright
  \setlength{\tabcolsep}{5pt}\renewcommand{\arraystretch}{1.3}\setlength{\arrayrulewidth}{0.5pt}
  \noindent\begin{tabular*}{\textwidth}{@{\extracolsep{\fill}}|c|c|c|c|}
    \hline
    \textbf{Coluna 1} & \textbf{Coluna 2} & \textbf{Coluna 3} & \textbf{Coluna 4} \\ \hline
    xxxxx & xxxxx & xxxxx & xxxxx \\ \hline
    xxxxx & xxxxx & xxxxx & xxxxx \\ \hline
    xxxxx & xxxxx & xxxxx & xxxxx \\ \hline
    xxxxx & xxxxx & xxxxx & xxxxx \\ \hline
    xxxxx & xxxxx & xxxxx & xxxxx \\ \hline
    xxxxx & xxxxx & xxxxx & xxxxx \\ \hline
    xxxxx & xxxxx & xxxxx & xxxxx \\ \hline
  \end{tabular*}\par
  \fonte{xxxxx (ano)}
\end{table}
